\section{从能运行的原型走向可演进的系统}

v1.0 已经把复位、装载、异常、分页、中断、Ring 3、调度、IPC、持久文件和
交互式 Shell 连成一条真实路径。它回答了“能否完全不用现成固件与引导器，
让自己的用户程序在 x86-64 上运行”。第二周期要回答的问题不同：

\begin{quote}
当对象数量增加、程序来自磁盘、地址页可以延迟出现、多个执行流共享资源，
并且磁盘写入可能被断电打断时，系统还能否保持身份、所有权和失败边界清晰？
\end{quote}

早期原型常把两个概念压进同一份存储，因为这能缩短第一次闭环。四个 PCB
槽位看起来就是四个 PID；八个描述符槽看起来就是八个打开文件；装载一段 ELF
看起来就意味着已经为整段分配物理页。这些做法并非错误，但它们只在对象少、
生命周期一致、没有共享时成立。

\begin{longtable}{L{0.22\textwidth}L{0.31\textwidth}L{0.35\textwidth}}
\toprule
\textbf{早期等价关系} & \textbf{扩展后为何失效} & \textbf{v2 长期关系}\\
\midrule
\endfirsthead
\toprule
\textbf{早期等价关系} & \textbf{扩展后为何失效} & \textbf{v2 长期关系}\\
\midrule
\endhead
PID 等于 PCB 槽位 & 槽位复用会让陈旧身份指向新对象 & PID 标识 Process，对象位置可以变化\\
Process 等于调度现场 & 一个地址空间可能有多个独立执行流 & Process 共享资源，Thread 参与调度\\
fd 等于文件句柄 & dup 和 fork 后需要共享偏移而独立关闭 & fd 引用 FileDescription，后者引用 Vnode\\
虚拟区等于物理页 & 未访问页面不应立即占用 RAM & VMA 描述意图，PTE 描述驻留事实\\
文件等于 inode 布局 & 挂载、设备与多种文件系统不共享磁盘格式 & VFS 对象隔离命名语义与存储格式\\
脏标志等于持久提交 & 断电可能把多块更新切在任意位置 & journal 记录事务边界和重放顺序\\
\bottomrule
\end{longtable}

因此 v2 不是简单增加系统调用数量。它要把早期被合并的身份拆开，并使每一条
引用都能回答四个问题：谁创建它，谁可以共享它，什么时候最后释放，任一步
失败后怎样恢复到调用前状态。

\section{v2.0 的边界与三类机器配置}

v2.0 仍从项目自己的 128 KiB ROM 开始，由自研 Stage 1 完成长模式与 ELF64
Kernel 装载。QEMU TCG 只模拟 x86-64 CPU、内存、传统中断控制器、时钟、
键盘与 IDE。Kernel 初始化后从自研 rootfs 执行 \code{/sbin/init}，PID1
再启动 \code{/bin/sh} 并回收孤儿。Shell 从磁盘执行普通工具，通过
fork、exec、pipe、dup、wait、signal 和 TTY 组合它们。

\begin{pdfdiagram}
  \node[os state] (boot) {ROM / Stage 1\\自研启动链};
  \node[os block,right=of boot] (kernel) {Kernel\\资源与硬件};
  \node[os state,right=of kernel] (init) {\code{/sbin/init}\\PID1};
  \node[os block,below=of init] (shell) {\code{/bin/sh}\\作业控制};
  \node[os state,left=of shell] (tools) {\code{/bin/*}\\外部 ELF};
  \node[os block,left=of tools] (fs) {rootfs v2\\journal};
  \draw[os arrow] (boot) -- (kernel);
  \draw[os arrow] (kernel) -- (init);
  \draw[os arrow] (init) -- (shell);
  \draw[os arrow] (shell) -- (tools);
  \draw[os arrow] (tools) -- (fs);
\end{pdfdiagram}
\diagramnote{Kernel 只建立最初的 Process/Thread；用户空间的后续结构由 PID1 和 Shell 通过稳定 ABI 形成。}

\subsection{Linux 对齐指分层能力，不指源码规模}

现代 Linux 能覆盖多架构、多处理器、网络、容器和大量设备。教学内核若把
“对齐”理解为复制这些表面数量，主线会很快失去可解释性。本项目只选择直接
决定长期正确性的结构：

\begin{itemize}
  \item Process 与 Thread 分离，调度器选择 Thread；
  \item 地址空间、文件表、工作目录和信号状态通过显式引用共享；
  \item fd、FileDescription 与 Vnode 分属三层身份；
  \item VMA 与页表项分别表示地址意图和当前硬件映射；
  \item 页缓存、块请求、journal 事务和磁盘块分别拥有状态；
  \item 容量由分配器和运行时资源限制决定，不由固定数组形状决定。
\end{itemize}

SMP、网络、USB、AHCI/NVMe、动态链接和完整 POSIX 不进入 v2.0。这个边界
不是降低质量，而是让单核系统中的每个共享对象、阻塞点和故障恢复都能够真正
完成并被证明。

\subsection{内存配置必须对应不同证据}

把所有配置都要求成“完整系统”会产生错误优化：为了让 64 MiB 跑最大并发，
实现可能重新引入狭窄固定表；为了让 64 GiB 启动，测试又可能忽略低内存下的
分配失败。三种正式配置因此承担不同角色。

\begin{longtable}{L{0.18\textwidth}L{0.16\textwidth}L{0.48\textwidth}}
\toprule
\textbf{配置} & \textbf{RAM} & \textbf{提供的证据}\\
\midrule
bootstrap & 64 MiB & 复位、启动链、异常、基础分页和历史失败路径仍能工作\\
functional & 256 MiB & PID1、Shell、VM、线程、信号、TTY 和持久化的完整功能\\
capacity & 64 GiB & 管理全部 E820 RAM、访问 4 GiB 以上地址并运行规模压力\\
\bottomrule
\end{longtable}

64 GiB 是当前主容量规格，不是硬编码实现上限。实际受管上界仍取 E820、
\code{CPUID.80000008H} 的物理地址宽度、64 TiB direct-map 与运行时资源
限制的交集。

\subsection{容量数字是压力目标，不是数组长度}

\begin{longtable}{L{0.25\textwidth}L{0.55\textwidth}}
\toprule
\textbf{领域} & \textbf{v2.0 容量配置验收下限}\\
\midrule
执行实体 & 同时至少 512 个 Thread，其中至少 256 个 Process\\
线程共享 & 单 Process 至少 64 个 Thread，具有独立栈、TLS 与 TID\\
描述符 & 分块动态表，默认 soft limit 256，hard limit 至少 4096\\
管道 & 至少 1024 条 64 KiB 管道，缓冲页按需获得\\
路径与参数 & 路径 4096 字节、组件 255 字节、参数环境合计 128 KiB\\
存储 & 至少 1 GiB 稀疏镜像、256 MiB rootfs、64 MiB 单文件\\
用户环境 & 至少 32 个独立 ELF 工具和 16 级流水线\\
\bottomrule
\end{longtable}

内部标识、容量、偏移和时间优先使用固定宽度的大类型；硬件寄存器、页表项和
磁盘字段则服从真实协议。达到限制必须返回明确错误，不能覆盖旧对象、静默截断
或只完成一半操作。

\section{路线不是功能清单，而是依赖图}

旧路线把 VFS 与新磁盘格式放在同一个阶段，把 COW 放在 VMA 之前，也把
\code{SYSCALL/SYSRET} 推迟到所有调用都完成以后。这些顺序都能写出短期代码，
却会制造第二次重构。新路线把十三个小版本组织成五个依赖波次：

\begin{longtable}{L{0.17\textwidth}L{0.20\textwidth}L{0.43\textwidth}}
\toprule
\textbf{波次} & \textbf{版本} & \textbf{建立的稳定边界}\\
\midrule
A：资源与实体 & v1.1--v1.3 & 内存生命周期、Process/Thread、动态对象与 fd\\
B：命名与来源 & v1.4--v1.6 & VFS、rootfs v2、PID1、磁盘 exec 和 wait\\
C：内存与组合 & v1.7--v1.9 & VMA/按需分页、COW、Unix I/O 和外部 Shell\\
D：并发与控制 & v1.10--v1.11 & 用户线程、futex、信号、TTY 与作业控制\\
E：持久与收口 & v1.12--v1.13 & 异步块层、journal、ABI 冻结与加固\\
\bottomrule
\end{longtable}

\begin{pdfdiagram}
  \node[os state] (a) {A\\资源与执行实体};
  \node[os block,right=of a] (b) {B\\VFS 与磁盘程序};
  \node[os state,right=of b] (c) {C\\VM 与用户组合};
  \node[os block,below=of c] (d) {D\\线程与交互控制};
  \node[os state,left=of d] (e) {E\\持久化与冻结};
  \node[os block,left=of e] (release) {v2.0\\集成发布};
  \draw[os arrow] (a) -- (b);
  \draw[os arrow] (b) -- (c);
  \draw[os arrow] (c) -- (d);
  \draw[os arrow] (d) -- (e);
  \draw[os arrow] (e) -- (release);
\end{pdfdiagram}
\diagramnote{每个箭头都表示后一个波次依赖前一个波次已经通过失败路径和回收验证的对象关系。}

\section{波次 A：先建立可回收资源与真实执行实体}

\subsection{v1.1：分配成功不等于生命周期完成}

v1.1 的第一增量已经解除低 64 MiB 页帧限制。2-bit 帧元数据根据最高可用
RAM PFN 动态计算，在低端启动映射中选择合法位置；正式页表从
\code{0xFFFF888000000000} 建立高半区 direct-map。64 GiB QEMU 已经在
4 GiB 以上执行分配、两种 64 位模式写入、读回和释放。

后续 Vnode、VMA、Thread、Pipe 和缓存页不能依赖 64 KiB 单调堆。v1.1
因此还要交付可分裂合并的页分配器、可释放内核堆或对象缓存、动态 Ring 0
栈、中间页表回收和统一资源快照。

分配器最容易被忽略的性质是失败原子性。例如申请四页时若第三页失败，前两页
不能遗留；释放范围必须先完整验证，再改变任意状态。固定种子模型至少执行
十万次分配、释放、分裂、合并、耗尽和错误释放，每轮结束后空闲页集合必须与
参考模型完全一致。

\subsection{v1.2：Process 是资源容器，Thread 才是调度实体}

历史 Unix 最初把单个进程既当资源所有者又当执行流。线程出现后，一个地址
空间中可以有多条寄存器现场和栈；若仍只保存一个 PCB，就无法说明哪一部分
应共享，哪一部分应独立。

\begin{pdfdiagram}
  \node[os block] (process) {Process\\PID};
  \node[os state,below left=of process] (address) {AddressSpace\\VMA / PTE};
  \node[os state,below=of process] (files) {FileTable\\FsContext};
  \node[os state,below right=of process] (signals) {SignalState\\children};
  \node[os block,right=of process] (thread) {Thread\\TID};
  \node[os state,below=of thread] (private) {CPU context\\kernel/user stack\\TLS / mask};
  \draw[os arrow] (process) -- (address);
  \draw[os arrow] (process) -- (files);
  \draw[os arrow] (process) -- (signals);
  \draw[os arrow] (process) -- node[above]{owns many} (thread);
  \draw[os arrow] (thread) -- (private);
\end{pdfdiagram}
\diagramnote{Process 退出和 Thread 退出不是同一个事件；最后一个 Thread 消失后才开始释放共享资源。}

v1.2 先在内核中建立这一区分，即使用户线程尚未开放。调度器的 Ready、
Running、Blocked 和 Zombie 状态属于 Thread；父子关系、地址空间和打开
文件属于 Process。PID 与 TID 都是 \code{uint64\_t} 身份，不等于数组位置。

\code{SYSCALL/SYSRET} 也在本阶段建立骨架。该指令不会像特权级变化的中断门
一样自动读取 TSS.RSP0，所以入口必须保存用户 RSP，切换到当前 Thread 的
内核栈，再保存 RIP 与 RFLAGS。STAR、LSTAR、FMASK 和 canonical 返回地址
都与 Thread 现场直接相关。若等到 v1.13 才加入，信号现场和所有新系统调用
都要重新验证一次入口边界。

\subsection{v1.3：fd、打开状态和文件身份必须分层}

fd 是 Process 局部小整数。FileDescription 保存当前偏移与打开状态，Vnode
表示命名空间中的文件身份。两次独立 open 应产生不同 FileDescription；
dup 或 fork 复制 fd 引用时则应共享同一个 FileDescription，因此共享偏移，
但每个 fd 可以独立关闭。

\begin{lstlisting}[language={},caption={动态描述符的引用方向}]
Process::FileTable[fd]
    -> FileDescription { offset, status_flags, reference_count }
        -> Vnode or Pipe or Console
\end{lstlisting}

动态并不表示“每次都无限增长”。FileTable 按块扩展，并有独立 soft limit
与 hard limit。限制是策略，存储是机制；测试可以在 256 MiB 下调低 soft
limit，却不能切换回固定 64 槽实现。随机 open、duplicate、close 和 limit
变化要证明陈旧 fd 不会因槽位复用而访问新对象。

\section{波次 B：先稳定命名空间，再改变磁盘格式}

\subsection{v1.4：VFS 先运行在旧文件系统适配器上}

Unix 把“路径是什么”和“磁盘字节怎样排列”分开，才允许同一个 open/read
接口访问磁盘文件、设备和伪文件。VFS 的 Superblock、Mount、Vnode、
Dentry/Path 与 FileDescription 分别表示文件系统实例、挂载关系、对象身份、
解析结果和一次打开状态。

v1.4 不改变 v0.11 磁盘格式，而是写一个 legacy-fs 适配器。这样绝对路径、
相对路径、\code{.}、\code{..}、根目录夹取、cwd、rename 和 unlink 的错误
可以在内存后端与旧磁盘后端上交叉验证。若 VFS 和 rootfs v2 同时改变，测试
发现 \code{/a/b} 解析错误时无法判断问题来自目录遍历还是间接块布局。

路径缓存也必须有失效规则。负缓存若记录“名字不存在”，create 后就必须失效；
rename 同时改变源父目录与目标父目录；mount 则改变越过目录项后的对象身份。
缓存只能加速已经定义的语义，不能成为语义来源。

\subsection{v1.5：rootfs v2 只改变持久布局}

rootfs v2 使用新 magic/version、固定宽度小端字段、4 KiB 逻辑块和可扩展
索引，支持至少 1 GiB 稀疏磁盘、256 MiB 根文件系统和 64 MiB 单文件。
启动区域与文件系统区域各有唯一所有者，mkfs 和镜像工具不能覆盖 ROM、
Stage 1 或 Kernel 容器。

这一阶段仍采用同步写入与明确 flush，不急于增加 journal。独立 fsck 从原始
字节重建 bitmap、inode、目录可达性和块所有权，避免“由写入实现检查自己的
输出”。旧 magic、未知非零介质、重复块、索引循环、孤儿 inode 和越界块都
必须安全拒绝，不能通过自动格式化掩盖损坏。

\subsection{v1.6：PID1 与磁盘 exec 终结内嵌程序脚手架}

正常 Kernel 只创建最初的 Process/Thread，从 rootfs 读取
\code{/sbin/init}。PID1 启动 Shell，接收失去父进程的子进程，并对 Zombie
执行 wait/reap。Zombie 已不能运行，大部分资源已经释放，但退出状态和最小
身份必须保存到父进程成功收取一次。

exec 沿用 Kernel ELF 装载阶段的两遍原则：第一遍读取并验证全部程序头、
W\^X、地址范围、重叠、入口和参数；第二遍建立新 AddressSpace 与初始
\code{argc/argv/envp} 栈。只有新映像完全可运行后才替换旧映像。路径错误、
畸形 ELF、128 KiB 参数环境超限或内存不足时，调用者仍继续执行原程序。

\section{波次 C：先解释地址意图，再做共享复制}

\subsection{v1.7：VMA 必须先于 COW}

页表项是硬件转换缓存的输入，只能表达某个虚拟页此刻是否 present、可写、
可执行以及映射哪个物理帧。它不能完整表达“这个区间来自 ELF 文件”“允许
按需增长”“是文件私有映射”或“尚未访问的匿名内存”。这些意图属于 VMA。

\begin{pdfdiagram}
  \node[os block] (fault) {用户 \#PF\\地址、访问类型、错误码};
  \node[os state,right=of fault] (vma) {查找 VMA\\区间、权限、来源};
  \node[os block,right=of vma] (classify) {分类\\非法 / 匿名 / 文件 / 栈 / COW};
  \node[os state,below=of classify] (page) {取得页\\zero / page cache / copy};
  \node[os block,left=of page] (pte) {提交 PTE\\TLB 失效};
  \node[os state,left=of pte] (resume) {恢复用户\\或明确终止};
  \draw[os arrow] (fault) -- (vma);
  \draw[os arrow] (vma) -- (classify);
  \draw[os arrow] (classify) -- (page);
  \draw[os arrow] (page) -- (pte);
  \draw[os arrow] (pte) -- (resume);
\end{pdfdiagram}
\diagramnote{处理程序先依据 VMA 判断访问是否合法，再依据 PTE 判断当前驻留状态；两者不能反过来。}

v1.7 建立按需 ELF、匿名 mmap、munmap、brk、受控栈增长和文件页缓存。大型
BSS 在未触碰时不占物理页，文件尾页未被文件覆盖的部分必须清零。VMA 拆分、
合并和 unmap 用十万步随机区间模型验证；exec、exit 和 unmap 后，VMA、
页表页、缓存引用和物理页都回到基线。

\subsection{v1.8：COW 是 VMA、PTE 与引用计数的协作}

fork 创建子 Process 时复制 VMA 元数据和资源引用，把私有可写驻留页在父子
PTE 中都改为只读 COW，并增加物理页引用。写故障到达时，处理器只报告
present/write 等硬件事实；Kernel 还必须从 VMA 与软件 PTE 位确认这是合法
COW，而不是对只读代码的非法写入。

若物理页只有一个引用，可以直接恢复可写权限；若仍有多个引用，才分配新页、
复制、替换 PTE 并递减旧页引用。分配失败必须保持旧 PTE 和引用计数不变。
多线程 Process 调用 fork 时，v2 明确只在子 Process 复制调用 Thread，
避免复制其他 Thread 正在执行的现场。

\subsection{v1.9：Unix I/O 让程序通过资源图组合}

pipe、dup、close-on-exec 与描述符继承完成后，Shell 可以先创建管道，再
fork 子进程，把子进程的 fd 0/1 重接到端点，关闭不需要的引用，最后 exec
磁盘工具。EOF 只有在最后一个写端引用关闭后出现；若 Shell 遗留一个写端，
读者将永久等待。

外部 Shell 只保留 \code{cd}、\code{export} 和 \code{exit} 等必须改变自身
状态的内建命令，至少 32 个普通工具位于 \code{/bin}。16 级流水线故意让
管道容量小于输入总量，从而证明生产者、消费者、阻塞唤醒与调度可以持续推进，
而不是依靠一次性大缓冲。

\section{波次 D：共享地址空间中的并发与终端控制}

\subsection{v1.10：线程需要 TLS，也需要原子的 compare-and-block}

用户 Thread 共享 Process 的 AddressSpace 与 FileTable，但拥有独立用户栈、
Kernel 栈、寄存器现场、TID、TLS 和 signal mask。线程入口使用 trampoline，
入口函数返回后统一 ThreadExit；最后一个 Thread 退出才进入 Process 退出。

大多数 mutex 获取在用户态原子指令上成功，不应每次陷入内核。发生竞争时，
futex wait 读取一个对齐的 32 位用户 word：只有值仍等于调用者期望值时才
加入等待队列。比较和阻塞必须相对于 wake 原子，否则 wake 可能落在“已经
比较但尚未排队”的窗口中，形成永久丢失唤醒。

\begin{lstlisting}[language={},caption={futex wait 必须保持的原子关系}]
lock(futex_bucket)
if (copy_from_user(word) != expected)
    return WouldBlock
enqueue(current_thread)
block_and_unlock(futex_bucket)
\end{lstlisting}

单 Process 的 64 Thread、全系统 512 Thread、TLS 隔离、虚假唤醒和地址
失效都进入验收。热竞争路径只累加计数，不按每次 wait/wake 打印日志。

\subsection{v1.11：信号和 TTY 把异步硬件事件变成用户语义}

PIT 提供 tick，但 sleep 不应每个 tick 扫描所有 Thread。绝对 deadline
进入有序 timer queue，到期时只唤醒相关等待者。串口日志的来宾时间也来自
这一单调时钟；时钟初始化前的日志明确标为 boot phase，避免伪造时间零点。

信号需要 disposition、mask、pending 和严格验证的用户 signal frame。
进程组、session 与 controlling TTY 决定 Ctrl-C、Ctrl-Z 应到达哪个前台
任务。Shell 自身通常不应被前台 Ctrl-C 终止；后台任务也不能抢走规范输入。
sigreturn 必须拒绝非 canonical RIP/RSP、非法 CS/SS 和提升 IOPL 的 RFLAGS。

\section{波次 E：先稳定写集合，再引入崩溃恢复}

\subsection{v1.12：journal 位于稳定 VFS 与页缓存之后}

同步 ATA PIO 会让调用 Thread 在端口状态上轮询。IRQ14 请求队列把设备操作
表示为独立 BlockRequest，提交者阻塞后其他 Ready Thread 仍可运行。缓存页
经历 loading、clean、dirty、writeback 和 error 状态，状态迁移必须与 I/O
完成和对象引用一致。

journal 只保证明确事务中的元数据写集合。事务描述和数据先落盘，commit
record 与设备 flush 完成后才算提交。崩溃测试在具名写点终止 QEMU，再用同一
镜像重启；恢复结果只能是完整旧状态或完整新状态。至少一千个固定种子断电点
用于验证重放幂等、bitmap/inode/目录一致和未提交事务不可见。

\subsection{v1.13：冻结 ABI，比继续加功能更重要}

ABI v2 固定系统调用号、寄存器约定、错误码、结构版本、字段宽度和对齐。
\code{SYSCALL/SYSRET} 成为正式入口；\code{INT 0x80} 只保留具名兼容测试
或明确移除。所有用户指针先复制或固定，长度加法、乘法和半开区间都进行溢出
检查，任何用户输入都不能使 Kernel panic。

\code{/dev} 与最小只读 \code{/proc} 或等价接口提供 Process、Thread、内存、
VFS、块层、信号和系统调用汇总。可观测性不等于刷日志：IRQ、每页故障和每次
futex 不打印；阶段完成、异常、速率受限告警和最终资源快照才形成结构化日志。
256 MiB 功能 soak 与 64 GiB 容量 soak 必须证明对象数不会随时间单调增长。

\section{每个版本怎样达到软件工程意义上的完成}

同一机制需要由不同层次的证据夹逼：

\begin{enumerate}[label=\arabic*.]
  \item ADR 记录问题、选择、替代方案、状态机、所有权和失败原子性。
  \item 宿主单元测试验证与硬件无关的边界，参考模型提供独立判断来源。
  \item 模块集成和固定种子随机测试交错创建、失败、共享、关闭与回收。
  \item ELF、ABI、汇编入口、磁盘镜像或 journal 格式由产物测试读取真实字节。
  \item QEMU 在适用内存配置上执行正常、容量、低资源和故障注入场景。
  \item 结束快照比较页、堆、Thread、Process、fd、Vnode、Pipe 与请求计数。
  \item 需求、架构、模块文档、教材、网站、手机 PDF 和发布记录来自同一提交。
\end{enumerate}

随机测试失败必须报告 seed、迭代位置、被破坏的不变量和可复现命令。QEMU
测试必须有总截止时间，串口标记要区分成功完成、预期拒绝与意外 panic。没有
资源回到基线的证据，“场景输出正确”仍不能算完成。

\section{v2.0 只发布已经稳定的机制}

v2.0 不新增主要对象和状态机。最终发布在同一干净提交上完成：

\begin{enumerate}[label=\arabic*.]
  \item 64 MiB 验证自研启动链和基础异常，256 MiB 执行完整功能套件。
  \item 64 GiB 管理全部可用 RAM，并在 4 GiB 以上执行真实读写回收。
  \item rootfs v2 启动 PID1、Shell 和至少 32 个磁盘 ELF 工具。
  \item 16 级流水线、重定向、前后台任务和 Ctrl-C 形成交互闭环。
  \item fork/COW、按需分页、文件映射、64 Thread 和 futex 压力后无泄漏。
  \item 至少 256 个 Process、512 个 Thread、4096 fd 和 1024 条管道达到
        明确容量下限。
  \item journal 在断电注入后重放，fsck 独立证明文件系统一致。
  \item 全部测试、ABI、格式文档、教材、手机导出和公开网站可复现。
\end{enumerate}

这样的版本号表示一组已经保存的工程证据，而不是功能完成比例。十三个小版本
让每次只跨越一个主要状态边界；v2.0 则把这些边界组合成一个能够启动、交互、
并发、持久化并在失败后恢复的完整教学系统。
